% =====================================================================
% Methods subsection: rigor-index tiering methodology (Week-5 / v3)
% Paste into the Methods section of the main paper (Overleaf).
%
% CITATION KEYS:
%   NEW  (BibTeX provided at the bottom of this file -- paste into your .bib):
%        jenks1967, fisher1958, glaser1963, cizek2007, nardo2008
%   EXISTING (already cited in the paper -- rename to YOUR key):
%        ho2020  -> your SEDA/Reardon citation (Ho, 2020 in the literature review)
%
% Uses natbib (\citep/\citet). If the paper uses plain \cite, replace them.
% Verify page numbers/years in Google Scholar before final submission.
% =====================================================================

\subsection{Classifying the Rigor Index into Ordinal Tiers}\label{sec:tiering}

The academic rigor classifier produces, for each school, a continuous composite score: each
component feature is standardized, the components are combined under the designed weighting
scheme, and weight is reallocated proportionally across whichever components a school actually
reports---the standard composite-indicator construction \citep{nardo2008}. Converting this
continuous score into the five ordinal tiers (Below Average, Average, Demanding, Very Demanding,
Most Demanding) is a separate methodological choice, and we treat it as one rather than as an
arbitrary partition.

We set the tier boundaries using Jenks natural breaks \citep{jenks1967}, the one-dimensional case
of Fisher's optimal-grouping algorithm \citep{fisher1958} and computationally equivalent to
one-dimensional $k$-means clustering. Natural breaks place the class boundaries so as to minimize
within-tier variance and maximize between-tier variance---that is, at the natural gaps in the
score distribution. We prefer this to equal-frequency quantiles for a reason internal to our own
review: \citet{ho2020} caution against distinguishing schools ``whose performance differs only
slightly,'' and equal-frequency quantiles do exactly that, forcing a tier boundary between
near-identical schools solely to fill equally sized buckets. Because the classification literature
also holds that the choice of break method is itself consequential, we report both the
natural-breaks tiers and an equal-quantile alternative together with their level of agreement
(49\% of scored schools receive the same tier under the two schemes), rather than presenting
either as canonical.

The resulting tiers are \emph{norm-referenced}: a school's tier reflects its position within the
scored population rather than a fixed external standard \citep{glaser1963}. Under natural breaks,
``Most Demanding'' corresponds to a composite score of at least approximately $1.5$ standard
deviations above the population mean and comprises roughly 700 schools, in contrast to the top
20\% that an equal-quantile scheme would assign. A norm-referenced tier shifts if the reference
population changes---here the scored population covers approximately 64\% of schools and is skewed
toward the recruiting universe---whereas a criterion-referenced alternative that fixes thresholds
on the underlying features (for example, a minimum mean AP examination score) would remain stable
across populations but requires separately justified cut scores \citep{cizek2007}. We adopt the
norm-referenced scheme for this iteration and flag the criterion-referenced alternative as an
available design decision for the client.

As an external check, mean SAT---a measure not used in constructing the tiers---rises
monotonically across the five natural-breaks tiers, from $1{,}052$ in Below Average to $1{,}303$
in Most Demanding with no inversions, indicating that the internally derived boundaries track
real between-school differences.

% =====================================================================
% BibTeX entries for the NEW references (paste into your .bib file):
% =====================================================================
%
% @article{jenks1967,
%   author  = {Jenks, George F.},
%   title   = {The Data Model Concept in Statistical Mapping},
%   journal = {International Yearbook of Cartography},
%   year    = {1967},
%   volume  = {7},
%   pages   = {186--190}
% }
%
% @article{fisher1958,
%   author  = {Fisher, Walter D.},
%   title   = {On Grouping for Maximum Homogeneity},
%   journal = {Journal of the American Statistical Association},
%   year    = {1958},
%   volume  = {53},
%   number  = {284},
%   pages   = {789--798}
% }
%
% @article{glaser1963,
%   author  = {Glaser, Robert},
%   title   = {Instructional Technology and the Measurement of Learning Outcomes: Some Questions},
%   journal = {American Psychologist},
%   year    = {1963},
%   volume  = {18},
%   number  = {8},
%   pages   = {519--521}
% }
%
% @book{cizek2007,
%   author    = {Cizek, Gregory J. and Bunch, Michael B.},
%   title     = {Standard Setting: A Guide to Establishing and Evaluating Performance Standards on Tests},
%   publisher = {SAGE Publications},
%   address   = {Thousand Oaks, CA},
%   year      = {2007}
% }
%
% @book{nardo2008,
%   author    = {Nardo, Michela and Saisana, Michaela and Saltelli, Andrea and Tarantola, Stefano and Hoffmann, Anders and Giovannini, Enrico},
%   title     = {Handbook on Constructing Composite Indicators: Methodology and User Guide},
%   publisher = {OECD Publishing},
%   address   = {Paris},
%   year      = {2008}
% }
